\chapter{Controle de Versão com Git e GitHub}
\label{ap:git}

Todo projeto de software — e firmware não é exceção — muda ao longo do tempo: funções são adicionadas, bugs são corrigidos, ideias são tentadas e, às vezes, abandonadas. Sem uma ferramenta para gerenciar esse histórico, é fácil perder uma versão que funcionava, sobrescrever o trabalho de outra pessoa, ou simplesmente esquecer \emph{por que} uma mudança foi feita. Este apêndice apresenta o \textbf{Git}, a ferramenta de controle de versão mais usada do mundo, e o \textbf{GitHub}, o serviço mais popular para hospedar repositórios Git na nuvem — ambos tão fundamentais para desenvolvimento de software moderno (incluindo projetos em PlatformIO) quanto o próprio compilador.

\section{O que é controle de versão}

Um \textbf{sistema de controle de versão} registra o histórico de alterações de um conjunto de arquivos ao longo do tempo, permitindo:

\begin{itemize}
    \item Voltar a qualquer versão anterior do projeto, a qualquer momento;
    \item Entender \emph{quem} mudou \emph{o quê}, \emph{quando} e \emph{por quê} (através de mensagens de commit);
    \item Trabalhar em uma funcionalidade nova sem arriscar quebrar a versão que já funciona (usando \textbf{ramificações}, ou \emph{branches});
    \item Colaborar com outras pessoas no mesmo código, combinando mudanças de várias fontes de forma organizada.
\end{itemize}

\begin{nota}
\textbf{Git} é a ferramenta que roda no seu computador e gerencia esse histórico localmente — nada obriga o uso de nenhum serviço externo. \textbf{GitHub} é, separadamente, um serviço \emph{online} que hospeda repositórios Git na nuvem, adicionando funcionalidades de colaboração (revisão de código, rastreamento de problemas, integração contínua). Existem alternativas ao GitHub — GitLab, Bitbucket —, mas todas são construídas sobre o mesmo Git.
\end{nota}

\section{Instalação e configuração inicial}

O Git está disponível para Windows, macOS e Linux — no Linux, geralmente já vem instalado ou disponível pelo gerenciador de pacotes da distribuição (\code{sudo apt install git}, por exemplo). Após instalar, identifique-se para que seus commits sejam atribuídos corretamente:

\begin{codigobash}{Configuração inicial (executada uma única vez por computador)}
$ git config --global user.name "Seu Nome"
$ git config --global user.email "seu.email@exemplo.com"
\end{codigobash}

\section{Os três estados de um arquivo}

Entender o Git fica muito mais simples ao visualizar as três áreas por onde um arquivo passa:

\begin{figure}[h]
\centering
\begin{tikzpicture}[
    box/.style={draw=secundaria, very thick, rounded corners=2mm, fill=fundoclaro,
                minimum height=1.5cm, minimum width=2.6cm, align=center, font=\small},
    node distance=1.6cm and 1.6cm
]
\node[box] (wd) {Diretório de\\Trabalho};
\node[box, right=of wd] (stage) {Área de\\\emph{Staging}};
\node[box, right=of stage] (local) {Repositório\\Local};
\node[box, right=of local] (remote) {Repositório\\Remoto\\(GitHub)};

\draw[-{Latex[length=2.5mm]}, thick, primaria] (wd) -- (stage)
    node[midway, above, yshift=1mm, font=\scriptsize\ttfamily] {git add};
\draw[-{Latex[length=2.5mm]}, thick, primaria] (stage) -- (local)
    node[midway, above, yshift=1mm, font=\scriptsize\ttfamily] {git commit};
\draw[-{Latex[length=2.5mm]}, thick, primaria] (local) -- (remote)
    node[midway, above, yshift=1mm, font=\scriptsize\ttfamily] {git push};

\draw[-{Latex[length=2.5mm]}, thick, acento]
    (remote.south) to[out=250, in=290]
    node[midway, below, yshift=-1mm, font=\scriptsize\ttfamily] {git pull / clone} (local.south);
\end{tikzpicture}
\caption{O fluxo básico do Git: do diretório de trabalho ao repositório remoto, e de volta.}
\end{figure}

\begin{itemize}
    \item \textbf{Diretório de trabalho} — os arquivos como estão no seu disco agora, editáveis livremente;
    \item \textbf{Área de \emph{staging}} (ou \emph{índice}) — uma zona intermediária onde você escolhe \emph{exatamente} quais mudanças farão parte do próximo commit, mesmo que outros arquivos também tenham sido alterados;
    \item \textbf{Repositório} — o histórico permanente de \emph{commits}, cada um uma ``fotografia'' completa do projeto em um instante específico.
\end{itemize}

\section{Criando um repositório}

\begin{codigobash}{Iniciando um repositório Git em um projeto existente}
$ cd meu-projeto-platformio
$ git init
Initialized empty Git repository in .../meu-projeto-platformio/.git/
\end{codigobash}

\code{git init} cria uma subpasta oculta \code{.git/}, onde todo o histórico do projeto será armazenado — apagar essa pasta apaga o histórico inteiro (o restante dos arquivos do projeto não é afetado).

\section{O ciclo básico: status, add e commit}

\begin{codigobash}{Verificando o estado do repositório}
$ git status
On branch main
Changes not staged for commit:
  (use "git add <file>..." to update what will be committed)
        modified:   src/main.cpp

Untracked files:
  (use "git add <file>..." to include in what will be committed)
        include/sensor.h
\end{codigobash}

\code{git status} é, disparadamente, o comando mais usado no dia a dia — ele mostra quais arquivos foram modificados, quais são novos (\emph{untracked}) e quais já estão prontos para o próximo commit.

\begin{codigobash}{Adicionando mudanças à área de staging}
$ git add src/main.cpp          # adiciona um arquivo especifico
$ git add include/sensor.h
$ git add .                     # adiciona TODAS as mudancas do diretorio atual
\end{codigobash}

\begin{codigobash}{Registrando um commit}
$ git commit -m "Adiciona leitura do sensor de luminosidade via ADC"
[main a1b2c3d] Adiciona leitura do sensor de luminosidade via ADC
 2 files changed, 34 insertions(+), 2 deletions(-)
\end{codigobash}

\begin{dica}
Uma boa mensagem de commit descreve \emph{o que} mudou e, quando não for óbvio, \emph{por que} — no modo imperativo (``Adiciona'', ``Corrige'', ``Remove''), como se estivesse dando uma instrução ao repositório. \code{"fix"} ou \code{"mudancas"} são mensagens que, seis meses depois, não ajudam ninguém (nem você mesmo) a entender o histórico do projeto.
\end{dica}

\section{Consultando o histórico}

\begin{codigobash}{Visualizando o histórico de commits}
$ git log --oneline
a1b2c3d Adiciona leitura do sensor de luminosidade via ADC
7e8f9a0 Configura watchdog timer no loop principal
3c4d5e6 Estrutura inicial do projeto PlatformIO
\end{codigobash}

\begin{codigobash}{Vendo exatamente o que mudou}
$ git diff                # mudancas ainda nao adicionadas ao staging
$ git diff --staged       # mudancas ja adicionadas, prontas para commit
\end{codigobash}

\section{Ignorando arquivos: \code{.gitignore}}

Nem todo arquivo gerado durante o desenvolvimento deveria entrar no histórico do Git — artefatos de compilação, por exemplo, podem sempre ser recriados a partir do código-fonte, e apenas poluiriam o repositório. Um projeto PlatformIO tipicamente ignora:

\begin{codigoini}{.gitignore de um projeto PlatformIO}
.pio/               ; pasta de compilacao, recriada a cada build
.vscode/            ; configuracoes locais do editor
*.o                 ; arquivos objeto intermediarios
compile_commands.json
\end{codigoini}

\begin{atencao}
Arquivos listados no \code{.gitignore} \textbf{deixam de ser sugeridos} pelo \code{git status} e pelo \code{git add .} — mas um arquivo que \emph{já} foi adicionado ao repositório antes de entrar no \code{.gitignore} continua sendo rastreado normalmente. Configure o \code{.gitignore} desde o início do projeto, antes do primeiro commit, para evitar esse tipo de surpresa.
\end{atencao}

\begin{esp32box}
Além de artefatos de compilação, \textbf{nunca} commite credenciais — como as constantes \code{WIFI\_SSID} e \code{WIFI\_PASSWORD} usadas no \cref{cap:wifi} e no projeto do \cref{cap:projeto-final}. A prática recomendada é isolá-las em um arquivo separado (por exemplo, \code{include/credenciais.h}), incluído no \code{.gitignore}, e versionar apenas um \code{credenciais.h.exemplo} com valores fictícios — assim, qualquer pessoa que clone o projeto sabe exatamente quais variáveis precisa preencher, sem que a senha real jamais chegue a ser registrada no histórico do Git (e, pior, publicada caso o repositório vá para o GitHub).
\end{esp32box}

\section{Ramificações (\emph{branches})}

Uma \textbf{branch} é uma linha independente de desenvolvimento — por padrão, todo repositório começa com uma branch principal (tradicionalmente chamada \code{main}). Criar uma nova branch permite experimentar ou desenvolver uma funcionalidade sem afetar o código já estável:

\begin{codigobash}{Criando e alternando entre branches}
$ git branch                       # lista as branches existentes
$ git checkout -b feature/wifi     # cria e ja muda para a nova branch
$ # ... faca commits normalmente nesta branch ...

$ git checkout main                # volta para a branch principal
$ git merge feature/wifi           # traz as mudancas de volta para main
\end{codigobash}

\begin{nota}
Versões mais recentes do Git introduziram \code{git switch} e \code{git restore} como alternativas mais específicas a \code{git checkout} (que, historicamente, acumulou funções demais). \code{git switch -c feature/wifi} é equivalente ao \code{git checkout -b feature/wifi} do exemplo acima.
\end{nota}

Quando duas branches modificam as mesmas linhas de um mesmo arquivo de formas diferentes, o \code{git merge} não consegue decidir sozinho qual versão manter — um \textbf{conflito de merge}. O Git marca o trecho conflitante diretamente no arquivo, e cabe à pessoa desenvolvedora escolher (ou combinar) manualmente qual versão faz sentido, antes de finalizar o merge com um novo commit.

\section{Trabalhando com o GitHub}

\subsection{Criando um repositório remoto}

Depois de criar um repositório vazio pelo site do GitHub, conecte-o ao repositório local e envie o histórico existente:

\begin{codigobash}{Conectando um repositório local ao GitHub}
$ git remote add origin https://github.com/seu-usuario/meu-projeto.git
$ git branch -M main
$ git push -u origin main
\end{codigobash}

\code{origin} é apenas um apelido (o mais comum, por convenção) para o endereço do repositório remoto; \code{-u} associa a branch local \code{main} à branch remota correspondente, para que futuros \code{git push}/\code{git pull} sem argumentos já saibam para onde ir.

\subsection{Clonando um repositório existente}

Para obter uma cópia completa de um projeto já hospedado no GitHub — o próprio código desta apostila, por exemplo —, junto com todo o seu histórico:

\begin{codigobash}{Clonando um repositório}
$ git clone https://github.com/algum-usuario/algum-projeto.git
$ cd algum-projeto
\end{codigobash}

\subsection{Enviando e recebendo atualizações}

\begin{codigobash}{Sincronizando com o repositório remoto}
$ git push          # envia seus commits locais para o GitHub
$ git pull          # traz e integra commits novos do GitHub
\end{codigobash}

\section{Colaborando com Pull Requests}

Em projetos com mais de uma pessoa, a forma padrão de propor uma mudança no GitHub é o \textbf{Pull Request} (PR): em vez de enviar commits diretamente para a branch principal, o fluxo típico é:

\begin{enumerate}
    \item Criar uma branch nova para a mudança (\code{git checkout -b feature/nome-da-mudanca});
    \item Fazer commits normalmente nessa branch, e enviá-la ao GitHub (\code{git push -u origin feature/nome-da-mudanca});
    \item Abrir um \textbf{Pull Request} pelo site do GitHub, comparando a nova branch com \code{main} — descrevendo o que mudou e por quê;
    \item Outras pessoas do projeto revisam o código, deixam comentários, e solicitam ajustes se necessário;
    \item Uma vez aprovado, o Pull Request é \textbf{mesclado} (\emph{merged}) à branch principal, incorporando a mudança de forma definitiva e registrada.
\end{enumerate}

\begin{dica}
Esse fluxo — branch dedicada, revisão antes de mesclar — vale a pena mesmo trabalhando sozinho em um projeto: ele força uma pausa para revisar as próprias mudanças antes de torná-las permanentes, e mantém a branch \code{main} sempre em um estado que compila e funciona.
\end{dica}

\section{Referência rápida}

\begin{table}[h]
\centering
\begin{tabular}{@{}ll@{}}
\toprule
\textbf{Comando} & \textbf{Finalidade} \\
\midrule
\code{git init}                  & inicia um repositório Git na pasta atual \\
\code{git clone <url>}           & copia um repositório remoto para o computador local \\
\code{git status}                & mostra o estado atual dos arquivos \\
\code{git add <arquivo>}         & move mudanças para a área de \emph{staging} \\
\code{git commit -m "..."}       & registra as mudanças em \emph{staging} como um novo commit \\
\code{git log --oneline}         & lista o histórico de commits, de forma resumida \\
\code{git diff}                  & mostra as mudanças ainda não adicionadas ao \emph{staging} \\
\code{git branch}                & lista, cria ou remove branches \\
\code{git checkout -b <nome>}    & cria e alterna para uma nova branch \\
\code{git merge <branch>}        & incorpora as mudanças de outra branch na atual \\
\code{git remote add origin <url>} & associa um repositório remoto ao projeto local \\
\code{git push}                  & envia commits locais para o repositório remoto \\
\code{git pull}                  & traz e integra commits do repositório remoto \\
\bottomrule
\end{tabular}
\caption{Comandos Git essenciais para o dia a dia de um projeto PlatformIO.}
\end{table}
